% !TEX root = ../main.tex
\lecture{2}{2026-09-24}{Woche 2}

\subsection{Nachfrage}

\begin{definition}
    [Gesetz der Nachfrage] Bei steigenden Preisen, tendieren Konsumenten zu alternativ Produkten (Substitutionseffekt). Nachfrage reduziert mit absteigende Kaufkraft (Einkommenseffekt)  
\end{definition}
Durch Funktiondarstellung wird Menge/Dauer und Zahlungsbereitschaft quantifiziert. \footnote{Zeitlich unabhängig, abstrahiert von allgemeine Tendenzen der Kaufkraft der breiten Bevölkerung. Geld immer noch vorausgesetzt.}  
\subsubsection{Vertikale Interpretation}
Sei eine Menge gegeben erhält man beim Ablesen die Information zur Zahlungsbereitschaft des "\textit{letzten Käufers}".

\textbf{Beispiel Kannabis}
Ersetzung von Verfahren mit einfacher Busse. Wahrgenommene Kosten gesunken. "Politik hat falsch das Gesetz der Nachfrage nicht berücksichtigt"

\subsection{Angebot}

\begin{definition}
    [Gesetz des Angebotes] 
\end{definition}

Funktion stellt Produktionskosten über Menge dar, sodass \textit{theoretisch} für jede weitere Einheit eines Produkts mehr Kosten verursacht.

\subsection{Herleitung Marktgleichgewicht}

\begin{definition}
    [\textit{magische} Marktgleichgewicht] optimaler Preis und Menge bei Schnittpunkt $D(Q) = S(Q)$. Darin ist die Angebotene Menge einer Ware gleich der nachgefragten Menge \footnote{Erklärt nicht die Wertverhältnissen zwischen verschiedenartigen Waren oder die zeitliche Kostenentwicklung. Besagt die Preisen werden durch die subjektive Zahlungsbereitschaft und den ProdKosten bestimmt. Kein systematischer urspung von wert und preis}
\end{definition}

\subsection{Effizienz}
\begin{definition}
    [Paretto Effizienz] unkritische Normative Anforderung. Eine Allokation ist \textit{Paretto effizient} wenn mit einer Bewegung der Konsum und Produkttionsmengen und -preise keine Person vernachlässigt werden kann.
\end{definition}

\textbf{Heisst nicht Fairness oder Gerechtigkeit auch wenn effizient aus Marktsicht.}

\subsection{Ungleichgewicht}
Bei ungleichgewichtsmoment sollten sich die Akteure richtung gleichgewichtbewegen. Da bei einer Menge ausserhalb des Gleichgewichts bei 1 mit ProdKosten 5 und Zahlungsbereitschaft 10 

\subsubsection{Höchstpreis}
Durch einführung eines Preisdeckels $\implies$ Überschüssigenachfrage
\subsubsection{Mindestpreis}
Durch erhöhung des Mindestpreises $\implies$ überschussangebot

\subsection{Marktversagen}
Zu sagen dass MGG Parettoeffizient ist braucht man starke vorrausetzungen.
Nur di
\begin{itemize}
    \item Vollkommene Konkurrenz
    \item keine externen effekten
    \item informationssymmetrie
\end{itemize}

\subsubsection{Gründe für Marktversagen (ineffiziente Allokation)}
\begin{itemize}
    \item Unvollkomene Konkurrenz
    \item externen effekten
    \item Informationsasymmetrie (insidertrading)
\end{itemize}

\subsection{Produzentenrente}

\begin{definition}
    [Produzentenrente] Differenz der Produktionskosten zum Gleichgewichtspreis pro Menge 
\end{definition}


\subsection{Konsumenten Rente}
\begin{definition}
    [Konsumentenrente] Differenz der Zahlungsbereitschaft zum Gleichgewichtspreis. Geldmass vom Monetären Vorteil der Teilnahme am Markt.
    \begin{equation}
        KonsumenteRente_{gesamt} = \int_{x_0}^{X_p}Demand(x)dx - \int_{x_0}^{X_p} P(x) dx
    \end{equation}
    wobei P$(x)$ eine funktion in Form y = a ist.
\end{definition}

\subsection{Determinaten Nachfrage}
Faktoren die die Zahlungsbereitschaft pro Menge verändern.
\begin{itemize}
    \item Einkommen 
        \subitem Mit steigendem Einkommen wächst die Zahlungsbereitschaft oder nachgefragte Menge \textit{\textbf{normaler Güter}}
    \item Individueller Geschmack
    \item Preis anderer Waren
        \subitem Nachfrage einer Ware steigt direkt mit steigende Substitutenpreise
        \subitem Nachfrage einer Ware steigt umgekehrt zu sinkende Komplementenpreise
\end{itemize}

\subsection{Determinanten Angebot}

\begin{itemize}
    \item Technologie
        \subitem weil es zum Sinken der Prodkosten führt \footnote{Wieso sinken dann kosten mit Technologie? noch keine erklärung.}
    \item Produktionspreise (Material, Löhne, Zinse)
\end{itemize}

\subsection{Komparative Statik}
Analyse oder Vergleich der alten Gleichgewichtslage mit der neuen GGL nach einer Verschiebung der $S(x)$ und $D(x)$ Funktionen.

\subsection{\textit{Wohlfahrts}verlust}
Noch genauer anschauen.
\subsection{Steuerlast}
Je höher die \textit{Elastität} oder flacher die Nachfragekurve desto kleiner die Konsumenten steuerlast. Weil die Fläche zwischen Nachfrage kurve und Nettopreisgerade nach Steuerlast kleiner ist, als die Angebotskurve;

